\section{Contract-Based and Runtime Assurance by Protected Object}
\label{sec:taxonomy}

We organize the reviewed literature by its \emph{protected object}: the
property or artifact that an approach is designed to preserve.
\Cref{tab:taxonomy} summarizes the twelve approach classes; the classes overlap,
and one system can contribute to more than one of them.

\begingroup
\footnotesize
\setlength{\tabcolsep}{3pt}
\renewcommand{\arraystretch}{1.08}
\begin{longtable}{@{}P{0.15\textwidth}P{0.15\textwidth}P{0.21\textwidth}P{0.21\textwidth}P{\dimexpr0.28\textwidth-24pt\relax}@{}}
  \caption{Reviewed approaches grouped by what they protect}
  \label{tab:taxonomy}\\
  \toprule
  Approach class & Protected object & Mechanism & Reported evidence & Obligations left to adjacent components \\
  \midrule
  \endfirsthead
  \multicolumn{5}{@{}l}{\textit{Table \thetable{} continued}}\\
  \toprule
  Approach class & Protected object & Mechanism & Reported evidence & Obligations left to adjacent components \\
  \midrule
  \endhead
  \midrule
  \multicolumn{5}{r@{}}{\textit{Continued on next page}}\\
  \endfoot
  \bottomrule
  \endlastfoot
    \multicolumn{5}{@{}l}{\textit{Mission intent}} \\*
    Mission contracts \cite{mallozzi2020crome,mallozzi2023contextual,mallozzi2023repair}
      & Goals, contexts, and component libraries
      & Contract-based goal graphs, refinement, synthesis, context switching, and repair
      & Mission examples and synthesized controllers; no physical team evaluation
      & Physical validity, event-time state, implementation identity, and peer reconciliation \\
    Language-conditioned planning \cite{mandi2024roco,obata2025lipllm}
      & Instructions, roles, task dependencies, and plans
      & Dialogue with environment feedback; dependency graph with a linear program
      & Multi-robot benchmarks, simulation, and selected real-world execution
      & Calibrated uncertainty, timing, compatibility, and deployment assurance \\
    Coordinated-mission assurance \cite{chaki2013assurance,chaki2013probabilistic,porter2013reliability,stancliff2006missionreliability}
      & Mission phases, coordination, utility, and team reliability
      & Markov-chain model checking, probabilistic verification, and reliability estimation
      & In-progress robot validation, measured robot behavior, and repairable-team analysis
      & Interface-carried uncertainty, provenance, software identity, lifecycle state, and recovery evidence \\
    \addlinespace[2pt]
    \multicolumn{5}{@{}l}{\textit{Component behavior}} \\*
    Component contracts and monitors \cite{luckcuck2026compositional,ghaffari2024rosmonitoring,caldas2024rvft}
      & ROS node assumptions, guarantees, and message traces
      & First-order contracts, composition rules, and generated monitors; topic and service monitoring; testing guidelines
      & Inspection-rover and smaller examples; UAV case study; literature, repository, and practitioner studies
      & Team epoch, resource ownership, state age, and evidence lineage \\
    Distributed programming and design \cite{ghosh2020koord,harbin2023design}
      & Coordination programs, proof obligations, and team architecture
      & Executable semantics, modular verification, generated infrastructure, and simulation-based robustness analysis
      & Three Koord applications; three simulated ATLAS case studies
      & Runtime freshness, evidence lineage, heterogeneous implementation identity, and recovery authority \\
    \addlinespace[2pt]
    \multicolumn{5}{@{}l}{\textit{Dynamical interaction}} \\*
    Dynamical contracts \cite{sharf2021dynamical,chen2019invariance,liu2022stlcontracts}
      & Input/output trajectories, invariant sets, and temporal tasks
      & Satisfaction and refinement checks, compositional invariance, and decentralized synthesis
      & Two-vehicle, microgrid, and numerical simulations
      & Message age, lifecycle, implementation identity, and recovery semantics \\
    Assume--guarantee learning \cite{kazemi2024agrl}
      & Neighbor abstractions, labels, and local policies
      & Regular-language contracts with a lower bound on global satisfaction probability
      & Grid-world, room-temperature, and traffic-signal case studies
      & Physical and middleware validity of the assumed abstractions \\
    \addlinespace[2pt]
    \multicolumn{5}{@{}l}{\textit{Distributed state}} \\*
    Middleware and QoS \cite{macenski2022ros2,omg2015dds,maruyama2016ros2,kronauer2021latency}
      & Data streams, delivery status, and endpoint compatibility
      & QoS policies and offered/requested matching; end-to-end latency measurement across DDS implementations
      & Standard and architecture descriptions; configuration-specific benchmarks
      & Event-time validity, physical decision validity, evidence lineage, and recovery \\
    Communication-aware coordination \cite{gielis2022communications,silva2026probabilistic,kaul2012aoi}
      & Information freshness and intermittent connectivity
      & Age-of-information analysis; temporal communication constraints in stochastic multi-robot planning
      & Critical review; analytical results; planning simulations
      & Implementation identity, evidence lineage, and coordinated recovery \\
    \addlinespace[2pt]
    \multicolumn{5}{@{}l}{\textit{Runtime authority}} \\*
    Runtime assurance and safety filters \cite{sha2001simplex,desai2019soter,shivakumar2020soterros,ames2017cbf,wang2017barrier,garg2024safecontrol}
      & Safety properties, safe sets, and control authority
      & Controller switching, barrier-constrained inputs, and learned safe control with verification
      & Drone surveillance in simulation and hardware; simulated multi-robot surveillance; wheeled-robot experiments; survey
      & Semantic evidence, implementation compatibility, and upstream reporting of interventions \\
    Mission-level runtime verification \cite{tsai2026missionrta,kekatos2026swarmrta}
      & Mission feasibility and swarm-level policies
      & Pre-execution command rejection; compositional verdicts with an explicit unknown
      & Simulated driving fault scenarios; simulated swarm mission
      & Physical feasibility, implementation identity, and recovery orchestration \\
    \addlinespace[2pt]
    \multicolumn{5}{@{}l}{\textit{Embodied execution}} \\*
    Cross-layer guardrails and runtimes \cite{kim2026guardrails,liu2026phyagentos,kim2026safeembodied}
      & Layered interventions, execution sessions, and evidence
      & Monitoring and intervention layers; compatibility preflight, session verification, and heartbeat
      & Position paper; simulated and physical embodiments; structured review
      & Common composition semantics across peer-robot interfaces \\
\end{longtable}
\endgroup


\subsection{Related Surveys}

Recent surveys use other organizing lenses.  Li \emph{et al.}\ classify
language-model research for multi-robot systems into task allocation, motion
planning, action generation, and human intervention \cite{li2026llmmrs}.
Street \emph{et al.}\ review formal models for robot teams under uncertainty,
and Gielis \emph{et al.}\ analyze how robotic applications and communication
systems interact \cite{street2023formal,gielis2022communications}.  Garg
\emph{et al.}\ organize learning-based safe control by safety notion, control
method, and verification tool \cite{garg2024safecontrol}.  Kim \emph{et al.}\
review safety for long-horizon manipulation by intervention locus (planning,
policy, and execution) and grade the strength of the supporting evidence
\cite{kim2026safeembodied}.  These surveys clarify capabilities, models,
communication, or intervention layers.  The present review uses a different
unit of comparison: the protected object, together with the information required
when its guarantee becomes another component's assumption.

\subsection{Mission Intent}

CROME captures mission requirements as contract-based goal graphs, checks
feasibility, refines goals from a library, and synthesizes controllers
\cite{mallozzi2020crome}; later work decomposes contextual missions into
context-specific sub-missions with correct-by-construction switching
\cite{mallozzi2023contextual}.  Contract operations can also refine a mission
against a component library or repair an unrealizable specification
\cite{mallozzi2023repair}.  These methods make goals, contexts, refinement, and
realizability explicit at the mission boundary; data age, provenance, and
implementation identity lie outside their scope.

Language-conditioned planners address an adjacent problem: turning open-ended
instructions into plans, role assignments, or optimization models.  RoCo
coordinates language-model agents through dialogue with environment feedback,
and LiP-LLM combines a language model with a predefined skill set, a dependency
graph, and a linear program \cite{mandi2024roco,obata2025lipllm}.  Both provide
capability evidence, while collision checking, skill definitions, and
feasibility solvers remain external.  Expressing intent is not the same as
certifying its physical or distributed execution.

Quantitative assurance of coordinated autonomy addresses the mission level from
another direction.  Chaki, Dolan, and Giampapa characterize coordination
phenomena and mission phases and evaluate discrete-time Markov-chain models by
probabilistic model checking \cite{chaki2013assurance}; related work computes
mission utility by probabilistic verification \cite{chaki2013probabilistic}.
Reliability studies estimate the performance of sequentially coordinated
missions from measured robot behavior and compare repairable team
configurations \cite{porter2013reliability,stancliff2006missionreliability}.
Within the reviewed corpus, this line is the closest mission-level predecessor
of the present question because it connects coordination structure to
quantitative mission evidence.  Its protected object is mission or team
reliability, whereas we ask which obligations must cross the heterogeneous
interfaces that realize such a mission.

\subsection{Component Behavior}

Luckcuck \emph{et al.}\ specify ROS nodes with first-order assume--guarantee
contracts, compose node contracts into system properties through inference
rules, and translate contracts written in the ROS Contract Language (RCL) into
executable runtime monitors \cite{luckcuck2026compositional}.  The approach
connects behavioral contracts to node implementations; cross-robot state age,
team membership, and heterogeneous deployment compatibility are not its central
objects.  ROSMonitoring~2.0 monitors topics and services, including message
order, with partial ROS~2 support \cite{ghaffari2024rosmonitoring}.  Caldas
\emph{et al.}\ contribute guidelines for runtime verification and field-based
testing of ROS-based systems that emphasize observability, instrumentation, and
reporting \cite{caldas2024rvft}.  Together, these works show that
specification, instrumentation, and evidence collection are separate
engineering concerns.

Distributed robot programming and model-driven design connect software
structure to team behavior more directly.  Koord provides platform-independent
abstractions for distributed shared memory, sensing, and control; its executable
semantics expose proof obligations, of which platform-specific ones are checked
with physics-based models and hybrid verification tools \cite{ghosh2020koord}.
ATLAS models team architecture, mission, and collective intelligence and uses
simulation to explore design choices and robustness \cite{harbin2023design}.
Both analyses remain tied to their languages, models, generated infrastructure,
and simulation settings; neither defines an evidence envelope for every
interface of a heterogeneous deployed stack.

\subsection{Dynamical Interaction}

Contracts for dynamical systems place assumptions on input trajectories and
guarantees on outputs.  Sharf \emph{et al.}\ verify satisfaction and refinement
of linear contracts by linear programming and prove safety in a two-vehicle
example \cite{sharf2021dynamical}.  Chen \emph{et al.}\ use parameterized
temporal-logic contracts to compute robust invariant sets compositionally for
network systems \cite{chen2019invariance}, and Liu \emph{et al.}\ synthesize
decentralized controllers for interconnected nonlinear systems from
signal-temporal-logic contracts \cite{liu2022stlcontracts}.  Assume--guarantee
reinforcement learning expresses contracts over neighboring components as
regular languages and combines local satisfaction probabilities into a lower
bound for the whole system \cite{kazemi2024agrl}.  These approaches address
interaction and decomposition directly, but they assume that the declared
neighbor abstractions, labels, communication, dynamics, and timing are available
as modeled.

\subsection{Distributed State}

ROS~2 provides modular robot communication through interchangeable middleware
implementations, most commonly DDS \cite{macenski2022ros2,omg2015dds}.  DDS
quality-of-service (QoS) policies govern reliability,
history, durability, liveliness, deadline, and lifespan.  These are
data-stream mechanisms; they do not establish application-level freshness or
mission validity.  The \emph{Deadline} policy concerns the expected period
between samples, not the age of the physical event a sample represents.  The
\emph{Lifespan} policy expires samples relative to their source time stamp,
which by default records the write call rather than the physical event; the
standard relies on sufficiently synchronized clocks and allows a reader to
substitute the reception time stamp when it detects that they are not
\cite{omg2015dds}.

ROS~2 latency studies show that process structure, executor behavior, and the
DDS implementation affect end-to-end timing
\cite{maruyama2016ros2,kronauer2021latency}.  Communication surveys find that
coordination and networking assumptions are often designed and evaluated
separately \cite{gielis2022communications}, and recent planning work builds
temporal communication constraints and intermittent connectivity into
stochastic multi-robot planning \cite{silva2026probabilistic}.  Age-of-information
research characterizes how fresh the newest available update is
\cite{kaul2012aoi}.  These works provide deployment semantics and measurements
rather than compositional robot contracts.  They nevertheless expose a necessary
bridge: a dynamical or safety guarantee that assumes current peer state needs an
application-level rule that relates event time, clock uncertainty, delivery, and
the physical decision.

\subsection{Runtime Authority}

The Simplex architecture uses a simple, verified controller to bound the
behavior of a more complex one \cite{sha2001simplex}.  SOTER makes this pattern
programmable for robotics: a decision module transfers authority from an
advanced controller to a safe controller when a monitor predicts a violation
\cite{desai2019soter}, and SOTER on ROS evaluates the pattern in simulated
multi-robot surveillance \cite{shivakumar2020soterros}.  Control barrier
functions and multi-robot barrier certificates instead constrain the admissible
inputs to keep a modeled safe set invariant \cite{ames2017cbf,wang2017barrier}.
These mechanisms give strong conditional safety guarantees, but the guarantees
depend on the dynamics model, state estimation, update rate, feasibility, and
timely peer state.  A local intervention can preserve one robot's safety while
invalidating the team plan unless it is reported upstream.

Mission-level runtime verification extends monitoring beyond immediate safety.
Tsai and Hariri reject driving commands that would make the mission infeasible
before they execute \cite{tsai2026missionrta}.  Kekatos \emph{et al.}\ compose
platform, squad, and mission monitors for robot swarms over a messaging fabric
that makes evidence loss observable, so that unsupported negative verdicts
become an explicit \emph{unknown} rather than an all-clear
\cite{kekatos2026swarmrta}.  This treatment of missing evidence is the closest
reviewed counterpart to the admission semantics of \cref{sec:admission}.

\subsection{Embodied Execution}

Kim \emph{et al.}\ argue for modular safety guardrails that monitor and
intervene across action, decision, and human-centered safety
\cite{kim2026guardrails}.  PhyAgentOS separates cognitive planning from physical
execution and provides compatibility preflight, supervised execution sessions,
evidence-grounded verification, layered safety, and heartbeat monitoring
\cite{liu2026phyagentos}.  It is the closest reviewed runtime architecture to
several obligations synthesized here.  PhyAgentOS, however, defines one
embodied-agent runtime, whereas the framework of \cref{sec:framework} is an
architecture-independent classification of obligations on both cross-layer and
peer-robot interfaces.  Within the reviewed works, we found no composition model
that jointly represents temporal validity, ownership, provenance, and recovery
across those interfaces.

\subsection{Synthesis}

The reviewed literature supplies mechanisms for mission realizability, node
behavior, dynamical composition, communication observation, and runtime safety.
Coverage becomes discontinuous where these mechanisms meet.  A mission guarantee
may end before physical feasibility; a dynamical guarantee may assume current
neighbor signals; a middleware status may omit the physical event time; and a
safety intervention may never return to task allocation.  \Cref{sec:gaps}
turns these discontinuities into five interface gaps.
